% ╔══════════════════════════════════════════════════════════════════════════╗
% ║   CORE/PREAMBLE.TEX — NÚCLEO DEL FRAMEWORK CampusTeX                     ║
% ║   Cargador de módulos. Compilar SIEMPRE con LuaLaTeX (dos pasadas).      ║
% ╚══════════════════════════════════════════════════════════════════════════╝
%
%  Cada módulo tiene una única responsabilidad y vive en core/.
%  El ORDEN de carga importa y está documentado línea a línea; no
%  reordenar sin entender la dependencia anotada.
%
%  RUTAS: \input resuelve rutas relativas al directorio de compilación
%  (courses/<curso>/), no a este archivo. Por eso los módulos se cargan
%  vía \CoreDir. Si un documento vive a otra profundidad, puede definir
%  \CoreDir antes de cargar este preámbulo:
%      \newcommand{\CoreDir}{../core}
% ══════════════════════════════════════════════════════════════════════════════
\providecommand{\CoreDir}{../../core}
\providecommand{\ConfigDir}{../../config}

% ── Validación temprana ────────────────────────────────────────────
% \NombreCurso es obligatorio y específico de cada curso: fallar aquí
% con mensaje claro evita errores crípticos en fancyhdr/hyperref.
\ifdefined\NombreCurso\else
	\errmessage{CampusTeX: define \string\NombreCurso\space en el main.tex
		ANTES de cargar core/preamble}
\fi

% ── Configuración global del proyecto ──────────────────────────────
% Academia, docente, ciclo y versión en UN solo lugar. Los valores del
% main.tex del curso (definidos antes de este punto) tienen precedencia.
\InputIfFileExists{\ConfigDir/project.tex}{}{%
	\errmessage{CampusTeX: no se encontro config/project.tex}}

% =====================================================================
%  core/layout.tex — Geometría de página y columnas
% =====================================================================
\usepackage[
a4paper,
top=2.5cm, bottom=2.5cm,
inner=3.0cm, outer=2.0cm,
headheight=14pt
]{geometry}

% Columnas — los ejercicios se maquetan a 2 columnas con regla sutil
\usepackage{multicol}
\setlength{\columnsep}{0.8cm}
\setlength{\columnseprule}{0.3pt}
\def\columnseprulecolor{\color{GrisMedio}}
      % geometry + multicol (independiente)
% =====================================================================
%  core/math.tex — Paquetes matemáticos
%
%  REGLA DE ORO con unicode-math + LuaLaTeX:
%    1. amsmath        (base)
%    2. mathtools      (antes de unicode-math, evita el Warning de orden)
%    3. cancel, systeme
%    4. unicode-math   (SIEMPRE AL FINAL — reemplaza completamente a amssymb)
%
%  NO cargar amssymb: genera 4 errores "Command already defined"
%  (\eth, \smallsetminus, \digamma, \backepsilon).
%
%  Con unicode-math los nombres de símbolos cambian:
%    amssymb \blacklozenge  →  unicode-math \mdlgblklozenge
%    amssymb \triangle      →  unicode-math \triangle  (este sí existe)
%    amssymb \blacktriangle →  unicode-math \blacktriangle (existe también)
% =====================================================================
\usepackage{amsmath}
\usepackage{mathtools}
\usepackage{cancel}
\usepackage{systeme}
\usepackage[math-style=ISO, bold-style=ISO]{unicode-math}

% --- Alias de compatibilidad amssymb → unicode-math ------------------------
% El contenido académico existente usa \square (amssymb); unicode-math
% lo llama \mdlgwhtsquare. Se define al inicio del documento porque
% unicode-math termina de montar su tabla de símbolos en \begin{document}.
\AtBeginDocument{%
	\ifdefined\square\else
	\newcommand{\square}{\mdlgwhtsquare}%
	\fi
}
        % amsmath → mathtools → unicode-math (AL FINAL de math)
% ══════════════════════════════════════════════════════════════════
%  core/fonts.tex — Selección de familias tipográficas
%
%  Valores por defecto del framework. Un curso puede sobreescribir
%  cualquiera definiendo el comando ANTES de \input del preámbulo:
%
%      % en courses/economia/main.tex, antes del preámbulo:
%      \newcommand{\FuenteSerif}        {EB Garamond}
%      \newcommand{\OpcionesFuenteSerif}{Numbers=OldStyle}
%
%  \providecommand solo define el valor si el curso no lo definió ya.
%
%  ORDEN: este módulo debe cargarse DESPUÉS de core/math.tex, porque
%  \setmathfont requiere unicode-math (que a su vez carga fontspec).
% ══════════════════════════════════════════════════════════════════
\usepackage{fontspec}

\providecommand{\FuenteSerif}        {Libertinus Serif}
\providecommand{\OpcionesFuenteSerif}{}
\providecommand{\FuenteMath}         {Libertinus Math}
\providecommand{\FuenteSans}         {Libertinus Sans}
\providecommand{\FuenteMono}         {InconsolataN}

\setmainfont{\FuenteSerif}[\OpcionesFuenteSerif]
\setmathfont{\FuenteMath}
\setsansfont{\FuenteSans}[Scale=MatchLowercase]
\setmonofont{\FuenteMono}[Extension=.otf,UprightFont=*-Regular,BoldFont=*-Bold,Scale=0.85]
       % selección de fuentes — REQUIERE unicode-math ya cargado
% ══════════════════════════════════════════════════════════════════
%  core/colors.tex — Paleta académica del framework (monocroma)
%
%  Identidad SOBRIA compartida con el Academic_Class_Framework: un único
%  acento «tinta pizarra» (casi negro) + grises tipográficos. Sin azules,
%  verdes, naranjas ni dorado: imprime idéntico en blanco y negro.
%
%  Se conservan los NOMBRES históricos (los usan cajas y cursos), pero todos
%  apuntan al acento o a un gris. ÚNICA fuente de verdad para colores; para
%  cambiar el acento, redefinir DESPUÉS del preámbulo.
% ══════════════════════════════════════════════════════════════════
\usepackage[dvipsnames, svgnames, table]{xcolor}

% Acento único y grises (paleta canónica) ---------------------------------
\definecolor{AzulPizarra}    {HTML}{1F2A37}   % ACENTO único (tinta pizarra)
\definecolor{AzulMedio}      {HTML}{5B5B5B}   % jerarquía secundaria (gris medio)
\definecolor{GrisCarbón}     {HTML}{1A1A1A}   % texto principal
\definecolor{GrisMedio}      {HTML}{B7B7B7}   % líneas y filetes suaves
% Fondos suaves: todos un mismo gris muy tenue -----------------------------
\definecolor{GrisClaro}      {HTML}{F7F7F7}   % fondo neutro de cajas
\definecolor{AzulClaro}      {HTML}{F7F7F7}   % (antes fondo de teoría)
\definecolor{Marfil}         {HTML}{F7F7F7}   % (antes fondo de ejemplos)
\definecolor{VerdeSuave}     {HTML}{F7F7F7}   % (antes fondo de tips)
\definecolor{NaranjaSuave}   {HTML}{F7F7F7}   % (antes fondo de advertencias)
% Antiguos acentos de color → el acento único (monocromo) ------------------
\definecolor{VerdeOliva}     {HTML}{1F2A37}   % (antes marco de tips)
\definecolor{NaranjaTierra}  {HTML}{1F2A37}   % (antes marco de advertencias)
\definecolor{DoradoAcadémico}{HTML}{1F2A37}   % (antes fórmulas)
\definecolor{RojoSuave}      {HTML}{1F2A37}   % (antes énfasis crítico)
      % xcolor + paleta — antes de todo lo que usa colores
% =====================================================================
%  core/boxes.tex — Sistema de cajas pedagógicas (tcolorbox)
%
%  Identidad SOBRIA (compartida con el Academic_Class_Framework): sin marcos
%  de color, sin arcos redondeados, sin «boxed title» ni sombras. Solo un
%  filete de acento a la izquierda + fondo gris muy tenue + rótulo en acento.
%  Todas las cajas comparten el mismo estilo; se distinguen por su título.
%
%  NOTAS (unicode-math): \blacklozenge→\mdlgblklozenge; \triangle y
%  \blacktriangle sí existen; usar \ensuremath{} para símbolos en el título.
% =====================================================================
\usepackage[most]{tcolorbox}
\tcbuselibrary{theorems, breakable, skins, listings}

% Estilo base único: filete de acento a la izquierda, fondo tenue, sin adornos.
\tcbset{campusbase/.style={
	enhanced, breakable, sharp corners, frame hidden,
	colback=GrisClaro, colbacktitle=GrisClaro, titlerule=0pt,
	borderline west={1.6pt}{0pt}{AzulPizarra},
	boxrule=0pt, left=10pt, right=8pt, top=6pt, bottom=6pt,
	before skip=8pt, after skip=8pt,
	coltitle=AzulPizarra, fonttitle=\bfseries\sffamily\small}}

\newtcolorbox{cajadef}[1][]{campusbase, title={Definici\'{o}n #1}}
\newtcolorbox{cajaprop}[1][]{campusbase, title={Propiedad #1}}
\newtcolorbox{cajaejemplo}[1][]{campusbase, title={Ejemplo #1}}
\newtcolorbox{cajatip}{campusbase, title={\ensuremath{\mdlgblklozenge}~Tip Preuniversitario}}
\newtcolorbox{cajaadvert}{campusbase, title={\ensuremath{\blacktriangle}~Error com\'{u}n}}
\newtcolorbox{cajarespuesta}{campusbase, title={Respuesta}}
\newtcolorbox{cajaformula}{campusbase, title={F\'{o}rmulas Clave}}

% Problema (sin título): mismo filete de acento.
\newenvironment{problema}{\begin{tcolorbox}[campusbase]}{\end{tcolorbox}}
       % cajas tcolorbox — usa la paleta
% ══════════════════════════════════════════════════════════════════
%  core/headers.tex — Encabezados y pies de página (fancyhdr)
%
%  Usa \NombreCurso y \NombreAcademia: deben estar definidos antes de
%  \begin{document} (los define el §B del main.tex y config/project.tex).
% ══════════════════════════════════════════════════════════════════
\usepackage{fancyhdr}
\pagestyle{fancy}
\fancyhf{}
\fancyhead[LE]    {\small\sffamily\color{AzulPizarra}\leftmark}
\fancyhead[RO]    {\small\sffamily\color{AzulPizarra}\rightmark}
\fancyhead[LO,RE] {\small\sffamily\color{GrisMedio}\NombreCurso\ --- \NombreAcademia}
\fancyfoot[LE,RO] {\small\sffamily\thepage}
\fancyfoot[C]     {\small\sffamily\color{GrisMedio}Material de uso interno}
\renewcommand{\headrulewidth}{0.4pt}
\renewcommand{\footrulewidth}{0.2pt}
\renewcommand{\headrule}{\color{AzulMedio}\hrule width\headwidth height\headrulewidth}
     % fancyhdr — usa \NombreCurso, \NombreAcademia
% ══════════════════════════════════════════════════════════════════
%  core/sections.tex — Formato de capítulos, secciones y subsecciones
% ══════════════════════════════════════════════════════════════════
\usepackage{titlesec}

% \chapter solo existe en book/report; la guarda permite usar el núcleo
% también con article (p. ej. plantilla de examen).
\ifcsname c@chapter\endcsname
\titleformat{\chapter}[display]
{\normalfont\Large\bfseries\sffamily\color{AzulPizarra}}
{\large\sffamily\color{AzulMedio}UNIDAD \thechapter}{6pt}
{\rule{\linewidth}{1.5pt}\vspace{4pt}\Huge}
[\vspace{2pt}\rule{\linewidth}{0.5pt}]
\fi

\titleformat{\section}
{\normalfont\large\bfseries\sffamily\color{AzulPizarra}}
{\color{AzulMedio}\thesection}{0.8em}{}
[\color{GrisMedio}\titlerule]

\titleformat{\subsection}
{\normalfont\normalsize\bfseries\sffamily\color{GrisCarbón}}
{\color{AzulMedio}\thesubsection}{0.6em}{}
    % titlesec — usa la paleta
% ══════════════════════════════════════════════════════════════════
%  core/lists.tex — Listas y alternativas de examen (A–E)
% ══════════════════════════════════════════════════════════════════
\usepackage{enumitem}
\setlist[itemize,1]{label=\textcolor{AzulMedio}{\textbullet},leftmargin=1.5em}
\setlist[enumerate,1]{label=\textcolor{AzulPizarra}{\arabic*.},leftmargin=1.8em}

% Lista de alternativas con formato exacto de examen de admisión
\newlist{alternativas}{enumerate}{1}
\setlist[alternativas]{label=\textbf{\Alph*)},leftmargin=2em,itemsep=2pt,parsep=0pt}
       % enumitem + alternativas A–E
% ══════════════════════════════════════════════════════════════════
%  core/tables.tex — Tablas profesionales
% ══════════════════════════════════════════════════════════════════
\usepackage{booktabs}
\usepackage{array}
\usepackage{tabularx}
\usepackage{multirow}
      % booktabs, tabularx, multirow
% =====================================================================
%  core/graphics.tex — Gráficos, TikZ y pgfplots
% =====================================================================
\usepackage{graphicx}
\usepackage{float}
\usepackage{tikz}
\usepackage{pgfplots}
\pgfplotsset{compat=1.18}
\usetikzlibrary{arrows.meta,shapes,positioning,calc,
	decorations.pathmorphing,backgrounds}

% Ruta estándar de imágenes del framework: los cursos pueden usar
% \includegraphics{logo} sin ruta si el archivo está en assets/images/.
\graphicspath{{../../assets/images/}{../../assets/logos/}{./}}
    % tikz, pgfplots, graphicx + \graphicspath a assets/
% =====================================================================
%  core/links.tex — Hipervínculos y metadatos PDF (hyperref)
%
%  ORDEN: hyperref se carga casi al final, siempre después de
%  unicode-math y titlesec.
%
%  Los metadatos PDF se generan desde los metadatos del curso: no es
%  necesario repetir \hypersetup en cada main.tex.
% =====================================================================
\usepackage[
colorlinks=true,
linkcolor=AzulPizarra,
citecolor=AzulMedio,
urlcolor=AzulMedio
]{hyperref}

\hypersetup{
	pdftitle  = {\NombreCurso\ --- \NombreAcademia},
	pdfauthor = {\NombreDocente},
	pdfsubject= {Material preuniversitario de \NombreCurso}
}
       % hyperref — SIEMPRE después de unicode-math y titlesec
% ══════════════════════════════════════════════════════════════════
%  core/typography.tex — Tipografía fina
% ══════════════════════════════════════════════════════════════════
\usepackage{microtype}
\usepackage{setspace}
\usepackage{parskip}
\usepackage{caption}
\setstretch{1.15}
  % microtype, setspace, parskip
% ══════════════════════════════════════════════════════════════════
%  core/macros.tex — Macros pedagógicas del framework
% ══════════════════════════════════════════════════════════════════

% ── Contador de ejercicios ─────────────────────────────────────────
% Se reinicia en cada capítulo (unidad). Uso: \ejercicio Enunciado…
% En clases sin \chapter (article, p. ej. exámenes) el contador es corrido.
\ifcsname c@chapter\endcsname
\newcounter{numejercicio}[chapter]
\else
\newcounter{numejercicio}
\fi
\newcommand{\ejercicio}{%
	\stepcounter{numejercicio}%
	\noindent\textbf{\sffamily\color{AzulPizarra}\thenumejercicio.}\quad}

% ── Portada de clase (TikZ) ────────────────────────────────────────
% Uso: \portadaclase{Tema}{Semana}{N.Unidad}{Nombre Unidad}
\newcommand{\portadaclase}[4]{%
	\newpage\thispagestyle{empty}
	\begin{tikzpicture}[remember picture,overlay]
		\fill[AzulPizarra]
		(current page.north west)
		rectangle ([yshift=-4.5cm]current page.north east);
		\fill[DoradoAcadémico]
		([yshift=-4.5cm]current page.north west)
		rectangle ([yshift=-4.8cm]current page.north east);
		\fill[GrisClaro]
		(current page.south west)
		rectangle ([yshift=2cm]current page.south east);
		\node[anchor=west,text=white]
		at ([xshift=2cm,yshift=-1.5cm]current page.north west)
		{\Large\sffamily\bfseries \NombreAcademia};
		\node[anchor=west,text=GrisMedio]
		at ([xshift=2cm,yshift=-2.5cm]current page.north west)
		{\normalsize\sffamily Curso: \textbf{\NombreCurso}
			\quad Docente: \textbf{\NombreDocente}};
		\node[anchor=west,text=GrisMedio]
		at ([xshift=2cm,yshift=-3.5cm]current page.north west)
		{\normalsize\sffamily Semana: \textbf{#2}
			\quad Ciclo: \textbf{\CicloActual}};
		\node[anchor=center,text=AzulPizarra]
		at ([yshift=-7cm]current page.north)
		{\Huge\sffamily\bfseries #1};
		\node[anchor=center,text=AzulMedio]
		at ([yshift=-8.2cm]current page.north)
		{\large\sffamily Unidad #3: #4};
	\end{tikzpicture}
	\vspace*{10cm}}
      % \ejercicio, \portadaclase — usa paleta y tikz

% ══════════════════════════════════════════════════════════════════════════════
% FIN DEL NÚCLEO
% ══════════════════════════════════════════════════════════════════════════════
